% Mathematical fragment. Requires amsmath, amssymb, amsthm and theorem environments.
% All constants below are fixed-character constants unless explicitly stated otherwise.

\section{Recovering a character from finitely many noisy positions}
\label{sec:noisy-grids}

The starting point is the symmetric Li--L\'evy process in the repository's
\path{Article/sections/li.tex}. Its exact-path decoder uses individual jumps
in progressively smaller size bands.  The theorem below replaces those jumps
by increments between finitely many observed positions.  It permits bounded
errors chosen after the entire path is seen.  The observation duration remains
fixed.  In particular, increasing the duration is not needed for consistency.

The arithmetic inputs are the signed zero-counting formula and the rational
prime-frequency formula proved in the repository, labels
\texttt{li:lem:bands} and \texttt{zer:lem:rational}.  Their application to the
Li coefficients uses the upstream Dirichlet-GRH theorem.  No simplicity,
spacing, linear independence, or disjointness assumption on zero ordinates is
made.  The probabilistic arguments below are independent of any formal
verification status of the upstream implementation.

\subsection{The process and an explicit estimator}

Let $\chi$ be a primitive Dirichlet character; the conductor-one character
represents $\zeta$.  Count all signed ordinates $\gamma$ of its nontrivial
zeros, including multiplicities.  Put
\[
 \theta(\gamma)=2\arctan\frac1{2\gamma}\quad(\gamma\ne0),\qquad
 \theta(0)=\pi,\qquad
 \nu_\chi=\frac12\sum_\gamma
       (\delta_{\theta(\gamma)}+\delta_{-\theta(\gamma)}).
\]
Thus a central zero of multiplicity $m$ contributes $m/2$ at each of
$\pi$ and $-\pi$.  Let $X=X^\chi$ be the real, symmetric, compensated
L\'evy process with no Gaussian component and
\[
 \mathbb E_\chi e^{iuX_s}
 =\exp\left[-s\int(1-\cos uy)\,\nu_\chi(dy)\right].
\]
Its small jumps have infinite total variation.  Their squared sizes are
integrable: $\int y^2\nu_\chi(dy)<\infty$.  The compensation is therefore
part of the definition.

Choose once and for all $w\in C_c^\infty(\mathbb R)$ supported on $[1,2]$,
$w\ge0$, with
$\int_1^2 w(v)\,dv=1$, and write $J_w=\int_1^2w(v)^2\,dv$.
For $0<|y|<\pi$ put
\[
 G(y)=\frac1{2\tan(|y|/2)},\qquad
 a_p=\frac{\pi\sqrt p}{\log p},\qquad
 f_{T,p}(y)=w\!\left(\frac{G(y)}T\right)
                    \cos(G(y)\log p),                         \tag{NG1}
\]
and set $f_{T,p}=0$ elsewhere.  Here $p$ is a prime and $T\ge2$ is a
chosen spectral height.  The compact support of $w$ makes $f_{T,p}$ smooth
on the whole real line.  A concrete alternative requiring only $C^4$
regularity for the first-order theorem is
\[
 w(v)=2772(v-1)^5(2-v)^5\mathbf1_{[1,2]}(v),\qquad
 J_w=2772^2\frac{(10!)^2}{21!}.                         \tag{NG2}
\]
Its derivatives through order four vanish at both endpoints.  For the
arbitrary-order construction later, use $C^\infty$ or the analogous
beta window with exponent at least $2m+3$ for correction order $m$.

Fix $t>0$.  At mesh $h=t/n$ observe positions $Y_0,\ldots,Y_n$ satisfying
\[
 \left|(Y_j-Y_{j-1})-(X_{jh}-X_{(j-1)h})\right|\le\eta
 \quad(1\le j\le n).                                      \tag{NG3}
\]
The errors can be random, dependent, or selected adversarially from the
whole path.  For example, position errors $|Y_j-X_{jh}|\le\varepsilon$
give (NG3) with $\eta=2\varepsilon$.  Define
\[
 \widehat D_{T,h,t}(p)
   =-\frac{a_p}{tT}\sum_{j=1}^n f_{T,p}(Y_j-Y_{j-1}),\qquad
 D^0_{T,h,t}(p)
   =-\frac{a_p}{tT}\sum_{j=1}^n f_{T,p}(X_{jh}-X_{(j-1)h}).   \tag{NG4}
\]
The estimator needs no conductor, zero table, zero multiplicities, or
candidate character.  It uses $n$ elementary evaluations per probe prime.
If each computed probe value has an additional absolute evaluation error
at most $\xi$, its effect on (NG4) is at most $a_p\xi/(hT)$.
Thus certified evaluation precision can be selected separately from the
observation-error bound; floating-point demonstrations are not an
automatic certificate at arbitrarily large $T$.

\begin{theorem}[Noisy finite-grid arithmetic recovery]
\label{thm:noisy-grid}
Fix $\chi$, $t>0$, $w$, and distinct primes $p_1,\ldots,p_k$.  All asymptotic
constants in this theorem may depend on these fixed choices.  Let
$h=t/n\to0$, $T\to\infty$, and suppose (NG3) holds.
\begin{enumerate}
\item The clean estimator has mean
\[
 \mathbb E_\chi D^0_{T,h,t}(p)
 =\Re\chi(p)+O\!\left(\frac{\log^2T}{T}
                          +hT^3\log^2T\right).               \tag{NG5}
\]
Whenever $hT^3\log^2T\le1$, its covariances satisfy
\[
 \operatorname{Cov}_\chi(D^0(p),D^0(r))
 =\mathbf1_{p=r}\frac{\pi pJ_w}{2t(\log p)^2}
                    \frac{\log T}{T}
  +O\!\left(\frac1{tT}+\frac{hT^2\log^2T}{t}\right).        \tag{NG6}
\]
\item Uniformly over every perturbation satisfying (NG3), if
$\eta\le1/(16T)$ then
\[
 \|\widehat D(p)-D^0(p)\|_{L^2}
 \le C\eta\left(T^2\log T+
                 T^{3/2}\sqrt{\frac{\log T}{t}}\right).     \tag{NG7}
\]
In particular, if
\[
 hT^3\log^2T\longrightarrow0,
 \qquad \eta T^2\log T\longrightarrow0,                     \tag{NG8}
\]
then all $\widehat D(p_j)\to\Re\chi(p_j)$ in mean square.
\item If the stronger conditions
\[
 hT^{7/2}(\log T)^{3/2}\longrightarrow0,
 \qquad \eta T^{5/2}\sqrt{\log T}\longrightarrow0            \tag{NG9}
\]
hold, then
\[
 \sqrt{\frac{tT}{\log T}}
 \bigl(\widehat D(p_j)-\Re\chi(p_j)\bigr)_{j=1}^k
 \quad\Longrightarrow\quad
 N_k\!\left(0,
  \operatorname{diag}_{j\le k}
       \frac{\pi p_jJ_w}{2(\log p_j)^2}\right).              \tag{NG10}
\]
Thus the limiting errors at distinct probe primes are independent.
\end{enumerate}
\end{theorem}

\begin{corollary}[Any vanishing error bound suffices]
\label{cor:any-noise}
For any $h\to0$ and any known $\eta\to0$, choose
\[
 T=\min\{h^{-1/4},\eta^{-1/3}\},                            \tag{NG11}
\]
where $\eta^{-1/3}=+\infty$ if $\eta=0$.  Then (NG9) holds and the
conclusions of the theorem follow.  There is no required relative rate
between the two vanishing quantities $h$ and $\eta$.
\end{corollary}

\begin{proof}
For large $n$, $T\to\infty$.  Since $T\le h^{-1/4}$,
$hT^{7/2}(\log T)^{3/2}\le
C h^{1/8}(\log(1/h))^{3/2}\to0$.
When $\eta>0$, the other inequality gives
$\eta T^{5/2}\sqrt{\log T}
\le C\eta^{1/6}\sqrt{\log(1/\eta)}\to0$.
\end{proof}

The result is consistent with the repository's finite-grid equivalence
theorem: every \emph{fixed} finite grid has overlapping statistical laws.
Here the number of samples tends to infinity, and the error probability
tends to zero.  No fixed finite grid is asserted to identify the character
with probability one.

\subsection{The two analytic estimates}

\begin{lemma}[Weighted arithmetic moments]
\label{lem:weighted-moments}
For fixed primes $p,r$, as $T\to\infty$,
\begin{align*}
 I_p(T):=\int f_{T,p}\,d\nu_\chi
   &=-\frac{T}{a_p}\Re\chi(p)+O(\log^2T),\\
 I_{pp}(T):=\int f_{T,p}^2\,d\nu_\chi
   &=\frac{T\log T}{2\pi}J_w+O(T),\\
 I_{pr}(T):=\int f_{T,p}f_{T,r}\,d\nu_\chi
   &=O(\log^2T)\qquad(p\ne r).                            \tag{NG12}
\end{align*}
Also, for $0<u\le1$,
\[
 \nu_\chi(|y|>u)\le C u^{-1}\log(2/u),\qquad
 \int_{|y|\le u}y^2\nu_\chi(dy)\le C u\log(2/u).           \tag{NG13}
\]
\end{lemma}

\begin{proof}
The signed zero count is
$N_\chi(v)=\frac v\pi\log v+C_\chi v+O_\chi(\log(v+2))$.
Since $G(\theta(\gamma))=|\gamma|$, the zero count and partial summation
give
\[
 \int w(G/T)^2\,d\nu_\chi
       =\frac{T\log T}\pi J_w+O(T).
\]
The rational prime-frequency formula in the repository gives, for each
fixed rational $x>1$ and $T\le v\le2T$,
\[
 \sum_{T<|\gamma|\le v}\cos(\gamma\log x)
 =-\frac{(v-T)\Lambda(x)\Re\chi(x)}{\pi\sqrt x}
       +O_{\chi,x}(\log^2T),                              \tag{NG14}
\]
where the sum includes multiplicity; the main term is zero unless $x$ is
an integral prime power.  Stieltjes integration by parts against $w(v/T)$
gives the first formula in (NG12), because $\int w=1$ and
$\int|w'|<\infty$.  Apply the same argument with $w^2$ and use
$\cos^2 z=(1+\cos2z)/2$ for the second formula.
For distinct primes, product-to-sum leaves frequencies $\log(pr)$ and
$|\log(p/r)|$; neither argument is an integral prime power, proving the
third.  Finally $\theta(\gamma)\asymp1/|\gamma|$ for large ordinates gives
the first bound in (NG13).  Integration of $2v\nu_\chi(|y|>v)$ from $0$
to $u$ gives the second.  Finite large atoms and central zeros are absorbed
in the constant.
\end{proof}

\begin{lemma}[A generator bound for moving probes]
\label{lem:generator}
Let $L$ be the generator of $X$, so
\[
 Lg(x)=\frac12\int
       [g(x+y)+g(x-y)-2g(x)]\,\nu_\chi(dy).                \tag{NG15}
\]
For every integer $r\ge1$, if a bounded smooth function $g$ satisfies
$\|g^{(j)}\|_\infty\le M A^j$ for $0\le j\le2r$, $A\ge2$, then
\[
 \|L^r g\|_\infty\le C_r M[A\log(2A)]^r.                \tag{NG16}
\]
For fixed $p,w,r$, both $f_{T,p}$ and any fixed finite product of such
probes satisfy the derivative hypothesis with $A=T^2$ and constant $M$.
Consequently,
\[
 \|L^r f_{T,p}\|_\infty\ll T^{2r}(\log T)^r.             \tag{NG17}
\]
\end{lemma}

\begin{proof}
Write $\Delta_y g(x)=g(x+y)+g(x-y)-2g(x)$.  For a set of $r$ increments,
apply the integral second-derivative representation to those $y_i$ with
$A|y_i|\le1$, and the bound $\|\Delta_y\|\le4$ to all the others.
Commutation of differences and derivatives gives
\[
 \|\Delta_{y_1}\cdots\Delta_{y_r}g\|_\infty
 \le4^r M\prod_{i=1}^r\min\{1,A^2y_i^2\}.
\]
By (NG13),
\[
 B(A):=\int\min\{1,A^2y^2\}\,\nu_\chi(dy)
 \le\nu_\chi(|y|>A^{-1})+
      A^2\int_{|y|\le A^{-1}}y^2\nu_\chi(dy)
 \ll A\log(2A).
\]
The repeated integral is absolutely convergent, proving (NG16).
On the support of $f_{T,p}$, $|y|\asymp T^{-1}$ and
$|G^{(j)}(y)|\ll_j T^{j+1}$.  Differentiating (NG1) therefore gives
$\|f_{T,p}^{(j)}\|_\infty\ll_{j,p,w}T^{2j}$.  All endpoint
derivatives match zero because the window is smooth.  Leibniz' rule treats
products.
\end{proof}

\subsection{Proof of the noisy-grid theorem}

Let $P_sg(x)=\mathbb E g(x+X_s)$.  Twice applying the semigroup fundamental
theorem gives
\[
 P_hg(0)=g(0)+hLg(0)+\int_0^h(h-s)P_sL^2g(0)\,ds.       \tag{NG18}
\]
For the smooth compact probes this identity follows directly first with
the L\'evy measure truncated away from zero; the bounds in
Lemma~\ref{lem:generator} justify passage to the limit.  In particular,
no density for the atomic measure $\nu_\chi$ is required.
Since $f_{T,p}$ vanishes in a neighbourhood of zero, $f_{T,p}(0)=0$ and
$Lf_{T,p}(0)=I_p(T)$.  The semigroup is a contraction, and (NG17) yields
\[
 \mathbb E f_{T,p}(X_h)
       =hI_p(T)+O(h^2T^4\log^2T).                        \tag{NG19}
\]
Multiplying by $-a_p n/(tT)$ proves (NG5).  Apply the same calculation to
$f_{T,p}f_{T,r}$ and use independent stationary increments.  Under
$hT^3\log^2T\le1$, (NG19) is $O(hT)$, so subtracting the product of the
means contributes $O(h/t)$ to the covariance of the estimators.
Equations (NG12) and (NG19) now give (NG6).

For the noise estimate put $r_T=1/(8T)$ and
\[
 N_T=\#\{j:|X_{jh}-X_{(j-1)h}|\ge r_T\}.
\]
The support of $f_{T,p}$ lies in $|y|\ge\theta(2T)\ge1/(4T)$.
If $\eta\le1/(16T)$ and the two evaluations of a probe at a clean and
perturbed increment differ, the clean increment therefore has magnitude
at least $r_T$.  The mean value theorem gives the pathwise bound
\[
 |\widehat D(p)-D^0(p)|
 \le\frac{a_p\eta\|f'_{T,p}\|_\infty}{tT}N_T
 \le\frac{C\eta T}{t}N_T.                                \tag{NG20}
\]
Split $X_h$ into jumps larger than $r/2$ and its compensated symmetric
small-jump part.  The former occurs with probability at most
$h\nu_\chi(|y|>r/2)$, while the latter has variance
$h\int_{|y|\le r/2}y^2\nu_\chi(dy)$.  Thus
\[
 \mathbb P(|X_h|\ge r)
 \le h\left(\nu_\chi(|y|>r/2)+
             r^{-2}\int_{|y|\le r/2}y^2\nu_\chi(dy)\right).
                                                                  \tag{NG21}
\]
At $r=r_T$, the parenthesis is $O(T\log T)$ by (NG13).  The count $N_T$
is binomial, independently of how the perturbations are selected, with
\[
 \mathbb E N_T\le CtT\log T,\qquad
 \mathbb E N_T^2\le C(tT\log T)^2+CtT\log T.
\]
Together with (NG20), this proves (NG7).  Equations (NG5)--(NG7) prove
mean-square consistency under (NG8).

For the central limit theorem, take any fixed linear combination of the
clean coordinates.  After multiplication by $\sqrt{tT/\log T}$ its
centred summands are bounded by $C/\sqrt{tT\log T}\to0$.
Their total variance tends, by (NG6), to the corresponding quadratic form
of the diagonal matrix in (NG10), because (NG9) implies
$hT^3\log T\to0$.  The bounded triangular-array Lindeberg theorem and
Cram\'er--Wold give the centred joint Gaussian limit.  Multiplying (NG5)
by $\sqrt{tT/\log T}$ shows that its bias vanishes under (NG9).
Finally (NG7) gives a vanishing $L^2$ perturbation after the same
normalization; Slutsky's theorem proves (NG10), uniformly for all errors
obeying (NG3).

\subsection{A finite-sample bound with specified input constants}

The following bound keeps its constants visible. It becomes a numerical error certificate only after upper bounds for these model-dependent quantities have been supplied. Upper bounds for its
input quantities can be substituted throughout; in a fixed finite
character family they can be replaced by family-wide maxima.

\begin{proposition}[Confidence bound including adversarial error]
\label{prop:grid-confidence}
Fix one prime $p$, let $f=f_{T,p}$, and put
\begin{gather*}
 M=\|f\|_\infty,\quad L_T=\|f'\|_\infty,
 \quad K_1=\|L^2f\|_\infty,\quad K_2=\|L^2(f^2)\|_\infty,\\
 I_2=\int f^2\,d\nu_\chi,\quad
 b_{\rm ar}=\left|-\frac{a_p}T\int f\,d\nu_\chi-\Re\chi(p)\right|,\quad
 V=t(I_2+hK_2/2),\\
 \beta_T=\nu_\chi(|y|>r_T/2)+
         r_T^{-2}\int_{|y|\le r_T/2}y^2\nu_\chi(dy),
 \qquad \mu=t\beta_T.
\end{gather*}
If (NG3) holds with $\eta\le1/(16T)$, then, for every $u>0$, with
probability at least $1-3e^{-u}$,
\begin{align*}
 |\widehat D(p)-\Re\chi(p)|\le {}&b_{\rm ar}
              +\frac{a_phK_1}{2T}
              +\frac{a_p}{tT}
                     \left(\sqrt{2Vu}+\frac43Mu\right)\\
 &+\frac{a_p\eta L_T}{tT}
                     \left(\mu+\sqrt{2\mu u}+\frac23u\right).
                                                               \tag{NG22}
\end{align*}
The event on which this holds can be chosen independently of the
perturbations, so the statement is simultaneous over all admissible
perturbations of a realized clean grid.
\end{proposition}

\begin{proof}
By (NG18), the clean bias is at most the first two terms.  Its unnormalized
centred sum consists of independent variables bounded by $2M$, with total
variance at most $V$.  Bernstein's inequality bounds that sum in absolute
value by $\sqrt{2Vu}+4Mu/3$ outside an event of probability $2e^{-u}$.
The binomial count in (NG20) has mean at most $\mu$ by (NG21).
The one-sided Bernstein inequality gives
$N_T\le\mu+\sqrt{2\mu u}+2u/3$ outside an event of probability $e^{-u}$.
Use (NG20) and the union bound.
\end{proof}

\subsection{Recovering the whole arithmetic label}

\begin{theorem}[Identification in any fixed finite candidate family]
\label{thm:finite-family}
Let $\mathcal C$ be a finite set of primitive Dirichlet characters, with
one representative from each conjugacy orbit.  There are finitely many
primes $p_1,\ldots,p_k$ for which the codewords
\[
 c_\chi=(\Re\chi(p_1),\ldots,\Re\chi(p_k)),\qquad\chi\in\mathcal C,
\]
are distinct.  Write $\Delta=\min_{\chi\ne\psi}\|c_\chi-c_\psi\|_\infty>0$,
with the singleton family interpreted trivially.  The nearest-codeword
rule
\[
 \widehat\chi=\operatorname*{argmin}_{\psi\in\mathcal C}
          \max_{j\le k}|\widehat D(p_j)-\Re\psi(p_j)|       \tag{NG23}
\]
with any fixed tie rule satisfies
\[
 \sup_{\chi\in\mathcal C}\sup_{\text{errors satisfying (NG3)}}
        \mathbb P_\chi(\widehat\chi\ne\chi)\longrightarrow0
                                                               \tag{NG24}
\]
under (NG8).  For a quantitative bound, apply (NG22) at all the probe
primes and require each right-hand side to be less than $\Delta/2$;
the error probability is then at most $3ke^{-u}$.
In particular, for all sufficiently large $T$ along any sequence obeying
(NG8), there is a constant $c>0$ such that the same uniform error
probability is at most
\[
             3k\exp\!\left(-\frac{ctT}{\log T}\right).     \tag{NG24a}
\]
The constant and the threshold are permitted to depend on the fixed
candidate family, the window, the probe primes, and the fixed duration.
\end{theorem}

\begin{proof}
If two primitive characters have the same real values at every prime,
induce them to the unit group modulo the least common multiple of their
conductors.  Every reduced residue class contains a prime, so their sums
with their conjugates coincide as functions on this finite group.
Linear independence of group characters implies that they belong to the
same conjugacy orbit; primitive antecedents are unique.  For each of the
finitely many distinct pairs, select one distinguishing prime.  Their
union gives the required finite set.  If the maximum estimation error is
less than $\Delta/2$, the correct codeword is uniquely nearest.
Theorem~\ref{thm:noisy-grid}, a finite union bound, and then the maximum of
finitely many constants prove (NG24).
For the exponential conclusion use $u=ctT/\log T$ in (NG22).
Uniformly over the finite family, $V=O(tT\log T)$ and
$\mu=O(tT\log T)$. The stochastic square-root term is $O(\sqrt c)$,
its linear term is $O(c/\log T)$, and the noise term is
$O(\eta T^2\log T)$. Both bias terms tend to zero by (NG8).
Choose $c$ sufficiently small compared with $\Delta^2$, then take $T$
large and apply the finite union bound.
\end{proof}

In particular, this applies to the entire finite family of primitive
characters of conductor at most a prescribed $Q$.  The family, time
duration, and prime probes are fixed in the asymptotic statements.  This
does not assert statistical constants uniform in a growing conductor
bound.  It resolves the fixed-family noisy-grid part of the repository's
question \texttt{li:q:paths}; its requests for optimal information and
uniform conductor/duration/probe-prime complexity remain separate.

\subsection{Higher-order mesh correction and its cost}

The leading discretization error can be cancelled by the standard
Richardson extrapolation identity.  The variance bound needed for the
result below uses the support of the probes, rather than an unjustified
Gaussian approximation at the accelerated meshes.

\begin{theorem}[Arbitrary fixed-order correction]
\label{thm:extrapolation}
Fix an integer $m\ge1$ and positive numbers $\alpha_1,\ldots,\alpha_m$
with $\sum\alpha_j=1$.  Split the observed interval into consecutive
disjoint blocks whose endpoints lie on the base grid, with lengths tending to $\alpha_jt$. In block $j$, use
nonoverlapping increments of length $jh$, discarding at most $jh$ at its
end, and normalize its estimator (NG4) by its actual used duration $t_j$.
Assume that the error in every used increment is at most $\eta$.
Set
\[
 b_j=(-1)^{j+1}\binom mj,\qquad
 \widehat D^{[m]}(p)=\sum_{j=1}^m b_j\widehat D_j(p).       \tag{NG25}
\]
For fixed $m,\chi,p,w,t,\alpha$, and $\eta\le1/(16T)$,
\begin{align*}
 \|\widehat D^{[m]}(p)-\Re\chi(p)\|_{L^2}
 \le C\bigg(&\sqrt{\frac{\log T}{T}}
          +\frac{\log^2T}{T}
          +h^mT^{2m+1}(\log T)^{m+1}
          +\eta T^2\log T\bigg).                         \tag{NG26}
\end{align*}
If the supplied data are noisy \emph{positions} with individual errors
bounded by $\varepsilon$, the hypothesis holds with $\eta=2\varepsilon$
for all the coarser increments as well.  If only elementary increment
errors are bounded by $\eta_0$, use $\eta=m\eta_0$ instead.
\end{theorem}

\begin{proof}
Repeated semigroup integration, justified by Lemma~\ref{lem:generator},
gives
\[
 \frac1s\mathbb E f_{T,p}(X_s)
 =\sum_{r=1}^m\frac{s^{r-1}}{r!}(L^rf_{T,p})(0)
       +O\bigl(s^mT^{2m+2}(\log T)^{m+1}\bigr).           \tag{NG27}
\]
The finite-difference identities
$\sum_{j=1}^m b_j=1$ and
$\sum_{j=1}^m b_jj^r=0$ for $1\le r<m$ cancel all the displayed
discretization terms except $Lf(0)=I_p(T)$.  Dividing by $T$ proves the
bias part of (NG26).

For variance, $|f_{T,p}(x)|\le M\mathbf1_{\{|x|\ge r_T\}}$ and (NG21)
give, for \emph{every} $s>0$,
\[
 \mathbb E f_{T,p}(X_s)^2\le M^2s\beta_T
                         \ll sT\log T.
\]
The variance of the normalized clean block estimator is therefore at
most $C\log T/(t_jT)$.  The disjoint blocks have independent increments,
so the variance of their weighted sum is at most
$C(\sum b_j^2/\alpha_j)\log T/(tT)$ for large $n$.
The pathwise perturbation argument (NG20), applied in each block and
then combined by the triangle inequality in $L^2$, proves the noise term.
Since $t,m,\alpha$ are fixed, its lower-order term
$\eta T^{3/2}\sqrt{\log T/t_j}$ is absorbed by
$C\eta T^2\log T$.  This establishes (NG26).
\end{proof}

\begin{corollary}[Rates approaching the quarter-power]
\label{cor:quarter-rate}
For every fixed $m$ and $0<\kappa<m/(4m+3)$, choose $T=h^{-2\kappa}$.
If $\eta=O(h^{5\kappa}/\sqrt{\log(1/h)})$, then
\[
 \|\widehat D^{[m]}(p)-\Re\chi(p)\|_{L^2}
       =O\!\left(h^\kappa\sqrt{\log(1/h)}\right).         \tag{NG28}
\]
Consequently every exponent $\kappa<1/4$ is achievable by a sufficiently
large \emph{fixed} correction order.  The constant is allowed to depend
on that order.  No claim is made with $m$ tending to infinity, or that
$1/4$ is an optimal statistical exponent.
\end{corollary}

\begin{proof}
The stochastic term in (NG26) is $O(h^\kappa\sqrt{\log(1/h)})$.
The arithmetic bias is smaller.  The grid bias divided by the stated
rate is
$O(h^{m-(4m+3)\kappa}(\log(1/h))^{m+1/2})\to0$.
The assumed noise bound treats the last term.
\end{proof}

For the common variance upper bound in the preceding proof, the optimal
block proportions and its amplification factor are explicit:
\[
 \alpha_j=\frac{|b_j|}{2^m-1},\qquad
 \min_{\alpha_j>0,\,\sum\alpha_j=1}
       \sum_j\frac{b_j^2}{\alpha_j}=(2^m-1)^2.             \tag{NG29}
\]
This follows from Cauchy--Schwarz and $\sum|b_j|=2^m-1$.
The factor is $1,9,49,225$ for $m=1,2,3,4$, respectively.  It quantifies
the cost of cancelling more mesh bias.  It minimizes this shared
variance \emph{bound}; a sharp Gaussian covariance at every accelerated
mesh has not been proved here.

\begin{corollary}[An explicit mesh--precision tradeoff]
\label{cor:balanced-risk}
Fix $m$.  For sufficiently small positive $h,\eta$, write
$\ell_h=\log(1/h)$ and $\ell_\eta=\log(1/\eta)$, and put
\[
 T_h=h^{-2m/(4m+3)}\ell_h^{-(2m+1)/(4m+3)},\qquad
 T_\eta=\eta^{-2/5}\ell_\eta^{-1/5}.
\]
Choose $T=\min(T_h,T_\eta)$; if $\eta=0$ omit the second choice.
Then the root-mean-square error in (NG26) is at most a constant times
\[
 \max\left\{
 h^{m/(4m+3)}\ell_h^{(3m+2)/(4m+3)},\quad
 \eta^{1/5}\ell_\eta^{3/5}\right\}.                       \tag{NG30}
\]
This is an achieved upper bound, with fixed-character constants; neither
the exponent $1/5$ for the error bound nor the limiting mesh exponent
$1/4$ is asserted to be minimax optimal.
\end{corollary}

\begin{proof}
At $T_h$ the stochastic term and the grid-bias term in (NG26) have the
order of the first quantity in (NG30).  At $T_\eta$ the stochastic term
and noise term have the order of the second.  For large $T$ the two bias
terms increase with $T$, while $\sqrt{\log T/T}$ decreases.  Taking the
smaller bandwidth bounds the total by the larger displayed rate.
The arithmetic bias $\log^2T/T$ is smaller than the stochastic term.
\end{proof}

\subsection{Relation to existing probability theory and remaining scope}

Small-time moment expansion and the use of smooth probes of a L\'evy
measure are established methods; see Figueroa-L\'opez,
\emph{Small-time moment asymptotics for L\'evy processes},
Statistics \& Probability Letters 78 (2008), 3355--3365,
DOI \texttt{10.1016/j.spl.2008.07.012}~\cite{Figueroa2008}. General high-frequency
L\'evy-measure inference is developed, for example, by Richard Nickl,
Markus Rei\ss, Jakob S\"ohl and Mathias Trabs in
\emph{High-frequency Donsker theorems for L\'evy measures},
Probability Theory and Related Fields 164 (2016), 61--108,
DOI \texttt{10.1007/s00440-014-0607-3}, arXiv:1310.2523~\cite{Nickl2016}.
Related reconstruction of continuous and jump components from high-frequency observations is treated by Gonz\'alez C\'azares and Ivanovs~\cite{Gonzalez2021}.
The contribution established here is their explicit combination with
the arithmetic moving probes (NG1): at a fixed observation duration,
finite noisy positions determine every selected Euler coefficient and
hence any fixed finite list of primitive-character labels, despite the
common leading small-time law.  The scale and error estimates are proved
for the atomic zero measures themselves.

The global generator bound in Lemma~\ref{lem:generator} is deliberately
used only where it proves the claimed assertion.  In particular,
Theorem~\ref{thm:extrapolation} supplies mean-square and classification
rates, not a sharp central limit theorem throughout the accelerated
range.  Exact one-band information, optimal sample/precision tradeoffs,
fixed nonvanishing observation noise, and conductor-uniform statistical
complexity remain open in this treatment.  No worldwide publication
priority or Lean formalization is claimed for the new theorems.
